\documentclass[10pt,a4paper]{amsart}
\usepackage[margin=19mm]{geometry}
\usepackage{amsmath,amssymb,mathtools}
\usepackage[scr=rsfs,cal=euler]{mathalfa}
\usepackage{xfrac}
\usepackage[unicode,hidelinks,bookmarks=false]{hyperref}
\newcommand{\ie}{i.e.\ }
\newcommand{\cf}{cf.\ }
\newcommand{\resp}{respectively\ }
\newcommand{\Subsection}{\subsection}
\providecommand{\nobreakdash}{\nobreak\hbox{-}}
\setlength{\emergencystretch}{2em}
\multlinegap0pt
\usepackage{bm}
\usepackage{bbm}
\usepackage{dsfont}
\usepackage{xspace}
\usepackage{cancel}
\usepackage{mathtools}
\usepackage[shortlabels]{enumitem}
\usepackage{amsthm}
\makeatletter
\@ifundefined{theorem}{\newtheorem{theorem}{Theorem}}{}
\@ifundefined{lemma}{\newtheorem{lemma}[theorem]{Lemma}}{}
\@ifundefined{proposition}{\newtheorem{proposition}[theorem]{Proposition}}{}
\makeatother
\newcommand{\Coll}{\operatorname{Coll}}
\newcommand{\Harm}{\operatorname{Harm}}
\newcommand{\LB}{\mathcal{L}_{\mathrm{B}}}
\newcommand{\Q}{\mathbb{Q}}
\newcommand{\R}{\mathbb{R}}
\newcommand{\aff}{\operatorname{aff}}
\newcommand{\interior}{\operatorname{int}}
\title[Elementary equivalence of convex bodies in affine and projec]{Elementary equivalence of convex bodies in affine and projective languages}
\author{David Victor Feldman}
\date{Source: arXiv:2607.25064v2 (CC BY 4.0)}
\begin{document}
\maketitle

\noindent\emph{Source and licence.} Elementary equivalence of convex bodies in affine and projective languages. Version arXiv:2607.25064v2 (CC BY 4.0), distributed under the Creative Commons Attribution 4.0 International licence, as recorded in the arXiv licence metadata for this exact version and shown on the abstract page https://arxiv.org/abs/2607.25064v2. Full rights notice: licenses/arxiv-2607.25064-CC-BY-4.0.txt.\par\smallskip

\noindent\textit{Editorial scope.} Counted unit: the Proposition whose source label is \texttt{prop:crq} in the article (source0.tex lines 1661--1671, source label \texttt{prop:crq}), delivered together with the article's own complete original proof of it (source0.tex lines 1673--1775, one \texttt{proof} environment of 103 lines). No mathematics is written, repaired or re-derived: statement and proof are the article's own text line for line. Required context retained uncounted: lem:harm-sound (lines 1411--1415); thm:harmonic (lines 1441--1448); lem:transport (lines 1551--1572); lem:net (lines 1608--1623). Every definition, display and label of those lines is retained unchanged. Omitted from this package and not redistributed: 1--1410, 1416--1440, 1449--1550, 1573--1607, 1624--1660, 1672--1672, 1776--3556. Nothing inside a retained excerpt is omitted or altered: every retained body is delivered exactly as the article's lines are written, so no omission receipt of any kind accompanies this package. Citations: the retained excerpts carry no citation command at all, so no bibliography entry of the article is transcribed and no reference is redistributed. Reference adaptation: none was needed and none was made; every reference inside the delivered unit resolves inside it, so no unresolved reference or citation is produced. Printed slips kept literal: no printed word, symbol, hypothesis or number of the counted statement or of its proof was altered. Layout and class: the article's own class is \texttt{amsart} with packages including amsmath, amssymb, amsthm, geometry, enumitem, hyperref; the wrapper is the lane's pinned \texttt{amsart} editorial wrapper at 10pt on A4, which restates the theorem-like environments the retained text needs and adds the article's own macro definitions that the retained text needs (\texttt{\textbackslash Coll}, \texttt{\textbackslash Harm}, \texttt{\textbackslash LB}, \texttt{\textbackslash Q}, \texttt{\textbackslash R}, \texttt{\textbackslash aff}, \texttt{\textbackslash interior}), with extra wrapper packages (bm, bbm, dsfont, xspace, cancel, mathtools, enumitem). The delivered numbering is the unit's own. Assets: no retained span contains a figure environment, a graphics-inclusion command, a PDF-page inclusion, a table or a TikZ picture; no image file is copied into this package, the package declares no asset dependency and no image file is redistributed with it. Licence: version arXiv:2607.25064v2 of this article is distributed under the Creative Commons Attribution 4.0 International licence, as recorded in the arXiv licence metadata for this exact version and shown on the abstract page https://arxiv.org/abs/2607.25064v2. A full rights notice is delivered at licenses/arxiv-2607.25064-CC-BY-4.0.txt. Characters: every retained excerpt is pure ASCII, so the delivered engine, XeLaTeX with its default Latin Modern text font, reports no missing character.
\begin{lemma}[Soundness]\label{lem:harm-sound}
Let $K\subseteq\R^n$ be convex and $a,b,c,d\in K$. If
$K\models\Harm(a,b;c,d)$ then $(a,b;c,d)=-1$; in particular $c\ne d$ and
$d\notin\{a,b\}$.
\end{lemma}

\begin{theorem}[Interior completeness]\label{thm:harmonic}
Let $n\ge2$, let $K\subseteq\R^n$ be convex with nonempty interior, and let
$a,b,c,d\in\interior K$ be collinear with $(a,b;c,d)=-1$. Then
$K\models\Harm(a,b;c,d)$, with quadrangle witnesses
$P,X,Y,Z\in\interior K$. Consequently, on quadruples of collinear
\emph{interior} points, the formula $\Harm$ defines exactly the harmonic
quadruples, uniformly in $K$ and in $n$.
\end{theorem}

\begin{lemma}[Interior perspectivity transport]
\label{lem:transport}
Let $n\ge2$, let $K\subseteq\R^n$ be convex with nonempty interior, let
$a,b,c,d\in\interior K$ be collinear and pairwise distinct on a line $\ell$,
and let $e\in\interior K\setminus\ell$. Write $\rho>0$ for a radius with
$\overline{B(e,\rho)}\subseteq\interior K$, and let $A=\aff(\ell\cup\{e\})$.
Then for every $\varepsilon\in(0,\rho)$, every prescribed linear order type
of the images, and every $R_0\ge1$ there exist a line $m\subseteq A$ not
through $e$ and points $a',b',c',d'\in B(e,\varepsilon)\cap m$ such that:
\begin{enumerate}[label=\textup{(\roman*)}]
\item $x'\in\ell(e,x)$ for each $x\in\{a,b,c,d\}$ \textup(perspectivity from
$e$\textup);
\item $(a',b';c',d')=(a,b;c,d)$;
\item the images realize the prescribed order along $m$, and the shape ratio
$|a'-b'|/|c'-b'|$ can be prescribed to lie in $[R_0,\infty)$, jointly with
\textup{(i)} and the containment in $B(e,\varepsilon)$.
\end{enumerate}
Moreover, for any points $e,m_1,m_2,x,x'$ of $K$ with $m_1\ne m_2$,
$\neg\Coll(m_1,m_2,e)$, $\Coll(e,x,x')$ and $\Coll(m_1,m_2,x')$, the point
$x'$ is the perspectivity image of $x$ from $e$ on the line $\ell(m_1,m_2)$,
so clauses of this form are sound.
\end{lemma}

\begin{lemma}[The harmonic net reaches every rational]
\label{lem:net}
Let $q\in\Q\setminus\{0,1\}$. There exist $N=N(q)\in\mathbb{N}$, a bound
$V=V(q)\in\mathbb{N}$, and a finite sequence of \emph{harmonic instructions}
$(i_j,k_j,l_j)_{j\le N}$ with indices in $\{-2,-1,0,1,\dots,j-1\}$, with the
following property. For any three distinct collinear points $p_\infty,p_0,p_1$
of $\R^n$, define $z_{-2}=p_\infty$, $z_{-1}=p_0$, $z_0=p_1$, and
recursively let $z_j$ be the harmonic conjugate of $z_{i_j}$ with respect to
$\{z_{k_j},z_{l_j}\}$, i.e., the unique point with
$(z_{k_j},z_{l_j};z_{i_j},z_j)=-1$. Then all $z_j$ are defined \textup(each
instruction's three inputs are pairwise distinct and none equals the conjugate
being taken at a fixed point of the involution\textup), the scale coordinate
of each $z_j$ in the chart $(p_\infty,p_0,p_1)\mapsto(\infty,0,1)$ is a
rational of absolute value at most $V$, and the final point $w:=z_N$ has
coordinate exactly $q$.
\end{lemma}

\begin{proposition}\label{prop:crq}
Let $n\ge2$ and let $K\subseteq\R^n$ be convex with nonempty interior. For
every rational $q\notin\{0,1\}$ there are $\LB$-formulas
$\mathrm{CR}_q(a,b,c,d)$ and $\mathrm{CR}_{>q}(a,b,c,d)$ such that, for all
collinear quadruples of pairwise distinct points of $\interior K$,
\[
K\models\mathrm{CR}_q(a,b,c,d)\iff(a,b;c,d)=q,
\qquad
K\models\mathrm{CR}_{>q}(a,b,c,d)\iff(a,b;c,d)>q .
\]
\end{proposition}

\begin{proof}
\emph{The formulas.} Fix $q$ and let $(i_j,k_j,l_j)_{j\le N}$, $V=V(q)$ be as
in Lemma~\ref{lem:net}. Both formulas have the shape
\[
\exists e\,\exists m_1 m_2\,\exists a'b'c'd'\,\exists z_1\cdots z_N\,
\bigl[\Pi\wedge\mathrm{Ord}\wedge\mathrm{Net}\wedge\Phi\bigr],
\]
where, writing $z_{-2}:=a'$, $z_{-1}:=b'$, $z_0:=c'$ and $w:=z_N$:
\begin{itemize}
\item $\Pi$ \textup(transport\textup): $\neg\Coll(a,b,e)$,
$m_1\ne m_2$, $\neg\Coll(m_1,m_2,e)$, and for each
$x\in\{a,b,c,d\}$ the clauses $\Coll(e,x,x')$ and $\Coll(m_1,m_2,x')$,
together with pairwise distinctness of $a',b',c',d'$;
\item $\mathrm{Ord}$: $B(b',c',a')$;
\item $\mathrm{Net}$: for each $j\le N$ the clause
$\Harm\bigl(z_{k_j},z_{l_j};z_{i_j},z_j\bigr)$;
\item $\Phi$ is $d'=w$ for $\mathrm{CR}_q$; and for $\mathrm{CR}_{>q}$,
\[
\Phi:\equiv\;d'\ne w\;\wedge\;
\begin{cases}
B(w,d',a')\leftrightarrow B(w,b',a'), & q<0,\\[1mm]
B(w,d',a'), & q>0 .
\end{cases}
\]
\end{itemize}

\emph{Soundness.} Suppose a witness exists. By the soundness clause of
Lemma~\ref{lem:transport}, $a',b',c',d'$ are the images of $a,b,c,d$ under
the perspectivity from $e$ onto $\ell(m_1,m_2)$; they are pairwise distinct
by the clauses, collinear, and
$(a',b';c',d')=(a,b;c,d)=:\chi$. By Lemma~\ref{lem:harm-sound}, each clause
of $\mathrm{Net}$ forces $(z_{k_j},z_{l_j};z_{i_j},z_j)=-1$, so by induction
and the uniqueness of harmonic conjugates each $z_j$ is exactly the $j$-th
point of the net of Lemma~\ref{lem:net} over the scale triple
$(a',b',c')$; in particular $w$ is the unique point of the line with
$(a',b';c',w)=q$. For $\mathrm{CR}_q$: $d'=w$ iff $\chi=q$, since
$x\mapsto(a',b';c',x)$ is injective.

For $\mathrm{CR}_{>q}$, coordinatize the line of $m$ affinely so that, using
$\mathrm{Ord}$, the positions are $b'=0<c'=\sigma<a'=\Delta$ for some
$0<\sigma<\Delta$. The chart position of scale coordinate $v$ is the
M\"obius function $\mu(v)=\Delta v/(v+\Delta/\sigma-1)$
\textup($\mu(0)=0$, $\mu(1)=\sigma$, $\mu(\infty)=\Delta$\textup), whose pole
is at $v^{*}=1-\Delta/\sigma<0$. The coordinate $\chi(u)$ of position $u$ is
the inverse M\"obius map, with pole at $u=\Delta$; it is strictly increasing
in $u$ on $u<\Delta$, with range $(v^*,+\infty)$ there, and strictly
increasing from $-\infty$ to $v^{*}$ on $u>\Delta$. We claim that in every
such configuration
\[
\{u:\chi(u)>q\}=
\begin{cases}
\text{the }\{w,a'\}\text{-arc of the line containing }b', & q<0,\\
\text{the open segment }(w,a'), & q>0 .
\end{cases}
\]
For $q>0$: since $v^{*}<0<q$, the threshold position
$\mu(q)=\Delta q/(q-v^{*})$ satisfies $0<\mu(q)<\Delta$, the far side
$u>\Delta$ carries only coordinates $<v^{*}<0<q$, and on the near side
monotonicity gives $\chi(u)>q\iff u\in(\mu(q),\Delta)$: the open segment
between $w$ and $a'$, which moreover never contains $b'$. For $q<0$ there are
two subcases. If $q>v^{*}$ then $w$ is at position $\mu(q)<0$ on the near
side; the far side again carries only coordinates $<v^{*}<q$, and
$\{\chi>q\}$ is the segment $(\mu(q),\Delta)$, which contains $b'$: the
$\{w,a'\}$-arc through $b'$. If $q<v^{*}$ then $\mu(q)>\Delta$, so $w$ lies
on the far side; now the entire near side has coordinates $>v^{*}>q$, and on
the far side $\chi(u)>q\iff u>\mu(q)$, so $\{\chi>q\}$ is the complement of
the closed segment $[a',w]$: again the $\{w,a'\}$-arc containing $b'$. Since
membership of a point in the segment between two collinear points is the
betweenness relation, the displayed sets are defined by exactly the two
clauses of $\Phi$ \textup(the biconditional with the reference point $b'$
expressing ``same $\{w,a'\}$-arc as $b'$''\textup), and $d'\ne w$ excises the
equality case. Hence $K\models\mathrm{CR}_{>q}(a,b,c,d)$ implies $\chi>q$,
and $K\models\mathrm{CR}_q(a,b,c,d)$ implies $\chi=q$. \textup(No clause
constrains $\sigma,\Delta$ beyond $\mathrm{Ord}$, and the dictionary above
was derived for arbitrary $0<\sigma<\Delta$, so soundness holds for every
witness.\textup)

\emph{Completeness.} Conversely, suppose $\chi=q$ \textup(respectively
$\chi>q$\textup); we must produce a witness. Choose $e\in\interior K$ off the
line of the quadruple, with $\overline{B(e,\rho)}\subseteq\interior K$, and
work in the plane $A$. Apply Lemma~\ref{lem:transport} with the order
$b',c',a'$, with shape ratio $\Delta/\sigma\in\{2V+2,\,2V+4\}$ chosen so
that $|\chi-(1-\Delta/\sigma)|\ge1$, and with $\varepsilon$ small; then
scale down further so that $\Delta\le\rho/(2+2|\chi|)$. The chord of $K$
on the transport line through $B(e,\varepsilon)$ contains the full window
$[-\Delta,\Delta]$ of positions around $b'$ deep inside $\interior K$. By the
choice of ratio, the pole satisfies $v^{*}\le-(2V+1)$, so every net value
$v\in[-V,V]$ of Lemma~\ref{lem:net} has position
$|\mu(v)|\le\Delta V/(V+1)<\Delta$: all net points $z_j$, including $w$
\textup(as $|q|\le V$\textup), lie in the window, hence in $\interior K$,
and each is an interior point of the chord. Each instruction of the net is a
harmonic quadruple of collinear interior points, so by
Theorem~\ref{thm:harmonic} the corresponding $\Harm$ clause holds with
interior quadrangle witnesses. The image $d'$ lies at position $\mu(\chi)$,
and $|\mu(\chi)|=\Delta|\chi|/|\chi-v^{*}|\le\Delta(1+|\chi|)\le\rho/2$
by the ratio choice, so $d'$ also lies on the chord, inside $K$; it lies in
$B(e,\varepsilon)$ automatically by Lemma~\ref{lem:transport}. Finally
$\Phi$ holds by the dictionary just established: if $\chi=q$ then
$d'=w$; if $\chi>q$ then $d'$ lies in the displayed set. All clauses are
satisfied, completing the proof. The formulas are uniform in $K$ and $n$,
as all constructions took place in the plane $A$ through the quadruple's
line and one auxiliary interior point.
\end{proof}

\end{document}
